\addcontentsline{toc}{chapter}{Abstract}\vspace{-1cm}
%Border
\begin{tikzpicture}[remember picture, overlay]
  \draw[line width = 4pt] ($(current page.north west) + (0.75in,-0.75in)$) rectangle ($(current page.south east) + (-0.75in,0.75in)$);
\end{tikzpicture}

This project presents a comprehensive, real-time fall detection system leveraging dual Texas Instruments IWR6843AOP FMCW mmWave radar sensors. Designed to overcome the occlusion and privacy limitations of traditional cameras, the system utilizes a robust track-based auto-calibration pipeline to fuse data from two independent radars—placed anywhere in a room—into a single unified 3D point cloud, without requiring manual measurement or hardware synchronization.

The system architecture features an optimized edge-computing pipeline that performs nearest-frame timestamp synchronization, calculates rigid spatial transforms using SVD and RANSAC, and merges the point clouds. To ensure high-quality data, redundant detections are merged via voxel downsampling (0.1 m grid), and multipath ghost points are eliminated using Statistical Outlier Removal (SOR). This self-healing pipeline automatically detects sensor drift and recalibrates in real time.

The unified point cloud is subsequently fed into a hybrid Transformer-CNN-LSTM deep learning architecture. This model extracts global temporal dependencies, local kinematic patterns, and sequential motion memory to classify sequences precisely as "Fall" or "No-Fall" events. Real-time predictions are pushed to a Supabase PostgreSQL database and visualized on a React dashboard via WebSocket channels. The end-to-end fusion and classification pipeline demonstrates sub-second latency, providing a highly scalable, privacy-first solution for continuous ambient safety monitoring.

\pagebreak